% !TEX root = ../../thesis.tex

\section{Preliminaries}\label{sec:p2-prelim}

Throughout this chapter, $(M^{n+1}, g)$ denotes a smooth, closed (compact without boundary) Riemannian manifold with $n \geq 2$. We use $\mathcal{Z}_{n}(M; \mathbb{Z}_2)$ to denote the space of mod~$2$ flat cycles in $M$, equipped with the flat norm $\mathcal{F}$.

\begin{definition}[Caccioppoli set]\label{def:p2-caccioppoli}
A measurable set $E \subset M$ is a \emph{Caccioppoli set} (set of finite perimeter) if
\[
\mathrm{Per}(E) := \sup\left\{ \int_M \chi_E \operatorname{div} \omega \,:\, \omega \in \mathfrak{X}(M),\; \|\omega\|_{C^0} \leq 1 \right\} < \infty,
\]
where $\chi_E$ denotes the indicator function of $E$. The \emph{perimeter of $E$ in an open set $U$}, denoted $\mathrm{Per}(E, U)$, is the total variation of $D\chi_E$ in $U$.
\end{definition}

By De~Giorgi's structure theorem, the distributional derivative $D\chi_E$ of a Caccioppoli set $E$ is given by $D\chi_E = \nu_E \, \mathcal{H}^n \mres \partial^* E$, where $\partial^* E$ is the \emph{reduced boundary} of $E$ (an $n$-rectifiable set), and $\nu_E$ is the outward unit normal defined $\mathcal{H}^n$-a.e.\ on $\partial^* E$. This allows us to identify $D\chi_E$ with an element of $\mathcal{Z}_n(M; \mathbb{Z}_2)$, which we denote by $\partial E := \nu_E \, \mathcal{H}^n \mres \partial^* E$. With this identification,
\[
\mathbf{M}(\partial E) = \mathrm{Per}(E) \qquad \text{and} \qquad \mathcal{F}(\partial E, \partial F) = \|\chi_E - \chi_F\|_{L^1} = \mathrm{Vol}(E \triangle F).
\]

\begin{definition}[Varifold]\label{def:p2-varifold}
An \emph{$n$-varifold} on $M$ is a Radon measure on the Grassmannian bundle $\mathrm{Gr}_n(M)$. We denote the space of $n$-varifolds by $\mathcal{V}_n(M)$. Given a cycle $\Gamma = \partial \Omega \in \mathcal{Z}_n(M; \mathbb{Z}_2)$, we write $|\Gamma|$ for the associated integral varifold:
\[
|\Gamma| = \mathcal{H}^n \mres \partial^* \Omega \otimes \delta_{T_x(\partial^* \Omega)}, \qquad \mu_\Gamma := \mathcal{H}^n \mres \partial^* \Omega.
\]
\end{definition}

We will switch freely between the equivalent notations $\mathbf{M}(\Phi(x))$ and $\mathrm{Per}(\Omega(x))$.

\begin{definition}[Stationary and stable varifolds]\label{def:p2-stationary}
A varifold $V \in \mathcal{V}_n(M)$ is \emph{stationary} if its first variation vanishes:
\[
\delta V(X) := \int_{\mathrm{Gr}_n(M)} \operatorname{div}_S X \, dV = 0 \quad \text{for every } X \in \mathfrak{X}(M).
\]
A stationary varifold $V$ is \emph{stable} if, in addition, the second variation $\delta^2 V(\varphi, \varphi) \geq 0$ for all compactly supported normal deformations $\varphi$ of $\mathrm{reg}\, V$.
\end{definition}

\begin{definition}[Tangent cone~{\cite[Chapter~8]{Sim84}}]\label{def:p2-tangent-cone}
Let $V$ be a stationary $n$-varifold in $M$ and $Z \in \spt\|V\|$. For $\rho > 0$, the \emph{blow-up} of $V$ at $Z$ with scale $\rho$ is the varifold $(\eta_{Z,\rho})_\# V$, where $\eta_{Z,\rho}(x) = (x - Z)/\rho$ in normal coordinates centred at $Z$. A \emph{tangent cone} of $V$ at $Z$ is any varifold limit $C = \lim_{j \to \infty} (\eta_{Z,\rho_j})_\# V$ for some sequence $\rho_j \to 0$. By the monotonicity formula, at least one such limit exists; it is a stationary cone in $\mathbb{R}^{n+1}$.
\end{definition}

\begin{definition}[Sweepout and width~{\cite[Definition~1.1]{CLS22}}]\label{def:p2-sweepout}
A \emph{sweepout} of $M$ is a map $\Phi \colon [0,1] \to \mathcal{Z}_n(M; \mathbb{Z}_2)$ continuous in $\mathcal{F}$-topology such that $\Phi(x) = \partial \Omega(x)$, where $\{\Omega(x) : x \in [0,1]\}$ is a family of Caccioppoli sets with $\Omega(0) = \varnothing$ and $\Omega(1) = M$. We denote by $\mathcal{S}$ the collection of all such sweepouts. The \emph{width} of $M$ is
\[
W = W(M) := \inf_{\Phi \in \mathcal{S}} \sup_{x \in [0,1]} \mathbf{M}(\Phi(x)).
\]
It is a consequence of the isoperimetric inequality that $W > 0$~\cite[Proposition~0.5]{DLT13}.
\end{definition}

\begin{definition}[Optimal nested volume parametrized (ONVP) sweepout~{\cite[Definition~1.2]{CLS22}}]\label{def:p2-onvp}
A sweepout $\{\Phi(x) = \partial \Omega(x)\}$ is called
\begin{itemize}
    \item \emph{optimal} if $\sup_{x \in [0,1]} \mathbf{M}(\Phi(x)) = W$;
    \item \emph{nested} if $\Omega(x_1) \subset \Omega(x_2)$ for all $0 \leq x_1 \leq x_2 \leq 1$;
    \item \emph{volume parametrized} if $\mathrm{Vol}(\Omega(x)) = x$ for every $x \in [0,1]$.
\end{itemize}
\end{definition}

The existence of ONVP sweepouts follows from the monotonisation technique of Chambers--Liokumovich~\cite{CL20} combined with an Arzel\`a--Ascoli compactness argument; see Theorem~\ref{thm:non-excessive-existence} below for the stronger statement that also guarantees the non-excessive property.

We recall the following criterion from de~Lellis--Tasnady~\cite{DLT13}, which gives conditions under which varifold convergence of Caccioppoli boundaries is guaranteed:

\begin{proposition}[Varifolds vs.\ Caccioppoli set limits~{\cite[Proposition~A.1]{DLT13}}]\label{prop:p2-DLT-integrality}
    Let $\{\Omega^k\}$ be a sequence of Caccioppoli sets and $U$ an open subset of $M$. Assume that
    \begin{enumerate}[label=\textup{(\roman*)}]
        \item $D\chi_{\Omega^k} \to D\chi_\Omega$ in the sense of measures in $U$;
        \item $\mathrm{Per}(\Omega^k, U) \to \mathrm{Per}(\Omega, U)$
    \end{enumerate}
    for some Caccioppoli set $\Omega$. Then the varifolds $|\partial^*\Omega^k|$ converge to $|\partial^*\Omega|$ in the sense of varifolds.
\end{proposition}
